\documentclass{article}
\pagenumbering{gobble}
\usepackage{enumerate}
\usepackage{xcolor}
\usepackage[margin=1.4in]{geometry}
\title{Statement on Diversity, Equity and Inclusion}
\author{Joel Villatoro}
\date{}

\usepackage{blindtext}
\usepackage{hyperref}
\begin{document}
\vspace{-0.5in}
\maketitle

I have always seen promoting equity and diversity as a vital component of my work. I am deeply committed to promoting an environment that is welcoming to individuals of all backgrounds. At your university, I will commit to furthering its mission to provide an equitable and inclusive environment for its faculty and students.

As a discipline, mathematics is faced with an ongoing crisis of diversity and equity. To be able to address this problem, it is vital to acknowledge the struggles faced by women and members of underrepresented ethnic and cultural groups. I recognize that people of many backgrounds often face serious challenges such as a lack of role models, implicit bias, and individual and systemic discrimination. 

There are a few different ways I try to promote inclusiveness and equity when teaching and mentoring students. One way is by recognizing how certain class formats can disproportionately advantage or disadvantage students based on their backgrounds. Generally speaking, traditional lectures result in worse outcomes for URM students. On the other hand, active learning has been shown to narrow achievement gaps for underrepresented students in mathematics\footnote{See \url{https://doi.org/10.1073/pnas.1916903117}}. 

In order to reap the benefits of more equitable formats, it is necessary for the instructor to be committed to providing an inclusive classroom environment. One way I try to accomplish this is by making sure that I talk to every student, rather than favoring certain students. I also try to show an active interest in my students' backgrounds and life experiences. I believe that students are much more likely to feel as if they belong if they have established a rapport with their colleagues. Therefore, I try to encourage student interactions both inside and outside the classroom. For example, I encourage students to form study groups and to work together when doing homework exercises. When splitting students into groups in class, I prefer to assign groups randomly to avoid the formation of cliques which can make students from underrepresented groups feel excluded.

In mathematics, I believe that maintaining a ``growth mindset'' is particularly important to establishing an inclusive classroom environment. To elaborate, women and students from certain underrepresented groups sometimes struggle with an internalized ``fixed mindset'' which tells them that their ability to succeed in mathematics is intrinsic and unchangeable. Although all students can experience these feelings, I have found that they are particularly pernicious among students from underrepresented groups. For example, I have seen this frequently with the high school students I have interacted with through the Puertas program. Many of them, especially young Latina women, have sadly internalized a belief that they are simply ``not good at math'' and that this is an immutable characteristic. It is important to confront this attitude head on and encourage students to instead adopt a growth mindset where they view mathematical ability as a skill that can be improved.

In my personal life, I benefited from having incredible role models. I also benefited from growing up in an environment that reinforced the idea that academic success was worthwhile and achievable. My mother moved to the United States from Mexico when I was five years old. She worked incredibly hard to provide for her four children as a single mother and we frequently had trouble making ends meet. Supported by student loans and government assistance, my mother enrolled in graduate school. Her pay as a graduate student was very low and it was not easy for her to support four children while also pursuing her education. However, she persevered and she eventually earned her PhD in microbiology.

Although our financial circumstances were poor, I was able to go to college and eventually get a PhD myself. Knowing what I know now, I am in awe of what she was able to accomplish. She is just one of many role models I have had that instilled in me the idea that I could succeed and that I could have a place in higher education. Having benefited from these things, I consider it an important responsibility to instill in others the idea that, they too are capable of success no matter their background.

At various times in life, I have tried to take active steps to promote equity and diversity. While I was an undergraduate student, I worked as a camp counselor for a day camp called Project Transformation. The goal of Project Transformation was to provide summertime childcare and educational services to parents from low-income communities around Oklahoma City. While a graduate student at the University of Illinois, I participated in a mentorship program for first year graduate students from underrepresented backgrounds. As an MPS-Ascend postdoc, promoting diversity was an important component of my grant application. While at Washington University I have worked with a program called Puertas that provides tutoring and mentorship for high schoolers from predominantly Latino communities. I served on the mathematics department's DEI committee and regularly attend the Field of Dreams Conference to support the university's efforts to recruit URM undergraduates. I also helped to restart the undergraduate math club which had become inactive since COVID. My primary motivation for this was to place myself in a role where I could steer the undergraduate mathematics community in a direction that was inclusive to people of all backgrounds.
\end{document}
